\documentclass[11pt,letterpaper]{article}
\usepackage[margin=0.75in]{geometry}
\usepackage{amsmath,amssymb}
\usepackage{booktabs}
\usepackage{hyperref}
\usepackage{xcolor}
\usepackage{enumitem}

\hypersetup{
    colorlinks=true,
    linkcolor=blue,
    filecolor=magenta,
    urlcolor=cyan,
}

\title{\textbf{START HERE}\\\
\large Westchester County Sidewalk Coverage Analysis\\\
Planning Decision Support Package}
\author{Arcanum Research Initiative}
\date{October 2025}

\begin{document}

\maketitle

\section*{Purpose of This Package}

This package provides \textbf{comprehensive decision-making tools} for assessing sidewalk adequacy in Metro-North Transit-Oriented Development (TOD) areas. Rather than a single number, you'll find spatial analysis, prioritization frameworks, investment scenarios, and comparative benchmarks to support planning decisions.

\vspace{1em}

\section*{Assessment Framework}

\subsection*{Coverage Analysis: Metro-North TOD Areas (0.5-mile radius)}

\begin{table}[h]
\centering
\begin{tabular}{@{}lrr@{}}
\toprule
\textbf{Coverage Status} & \textbf{Road Count} & \textbf{Percentage} \\
\midrule
\textcolor{red}{\textbf{No Coverage (Priority)}} & 502 roads & 45.1\% \\
\textcolor{orange}{\textbf{One-Side Coverage (Opportunity)}} & 352 roads & 31.5\% \\
\textcolor{green!70!black}{\textbf{Both-Sides Coverage (Adequate)}} & 263 roads & 23.5\% \\
\midrule
\textbf{Any Coverage} & 615 roads & 54.9\% \\
\textbf{Total TOD Roads Analyzed} & 1,117 roads & 100\% \\
\bottomrule
\end{tabular}
\end{table}

\subsection*{Key Planning Context}

\begin{itemize}[leftmargin=*]
    \item \textbf{Transit Connectivity Gap:} 502 TOD roads (45.1\%) lack pedestrian access to Metro-North stations
    \item \textbf{Quick Win Potential:} 352 roads already have one-side coverage - completing the second side is cost-effective
    \item \textbf{Comparative Performance:} TOD areas have 3.0x better coverage than non-TOD areas (54.9\% vs 18.2\%)
    \item \textbf{Industry Context:} Best practice for TOD pedestrian coverage is 75-85\% (current: 54.9\%)
    \item \textbf{Equity Concern:} 81.8\% of non-TOD roads lack sidewalk coverage
\end{itemize}

\subsection*{Assessment: \textcolor{orange}{\textbf{MODERATE ADEQUACY - Substantial Investment Needed}}}

While TOD areas outperform the county average, nearly half of transit-accessible roads lack sidewalks. This limits walkability, transit ridership potential, and equitable access.

\vspace{1em}

\section{Decision-Making Tools in This Package}

\subsection{Excel Files (Interactive Data Exploration)}

Located in \texttt{Excel/} folder - open with Microsoft Excel, Google Sheets, or LibreOffice Calc:

\begin{enumerate}[leftmargin=*]
    \item \texttt{1\_{EXECUTIVE}\_SUMMARY.xlsx}
    \begin{itemize}
        \item Complete coverage breakdown by TOD vs Non-TOD
        \item Filter and sort roads by coverage status
        \item Use for: Investment prioritization, budget justification
    \end{itemize}

    \item \texttt{2\_TOD\_COMPARISON.xlsx}
    \begin{itemize}
        \item Side-by-side TOD vs Non-TOD vs County-Wide analysis
        \item Includes industry benchmarks and gap-to-target calculations
        \item Use for: Strategic planning, stakeholder presentations
    \end{itemize}

    \item \texttt{3\_ROAD\_TYPE\_ANALYSIS.xlsx}
    \begin{itemize}
        \item Coverage breakdown by road classification (paved, unpaved, elevated, etc.)
        \item Identifies road types with lowest coverage (unpaved: 4.0\%)
        \item Use for: Targeted investment by infrastructure type
    \end{itemize}

    \item \texttt{4\_AREA\_ANALYSIS.xlsx}
    \begin{itemize}
        \item Coverage statistics by road area (acres)
        \item Useful for cost estimation based on square footage
        \item Use for: Budget planning, ROI analysis
    \end{itemize}
\end{enumerate}

\textbf{All Excel files follow the Druck standard: ONE SHEET per file for clarity.}

\subsection{PDF Reports (Analysis \& Recommendations)}

Located in \texttt{Reports/} folder - open with any PDF reader:

\begin{enumerate}[leftmargin=*]
    \item \texttt{Executive\_Summary.pdf} (5 pages)
    \begin{itemize}
        \item High-level overview of coverage assessment
        \item Key findings and immediate recommendations
        \item Use for: City Council presentations, public communications
    \end{itemize}

    \item \texttt{Technical\_Analysis.pdf} (12 pages)
    \begin{itemize}
        \item Detailed methodology (DVRPC sidewalk-to-road ratio)
        \item Statistical validation and quality assurance
        \item Coverage by road type, spatial patterns, policy recommendations
        \item Use for: Technical review, grant applications, methodology defense
    \end{itemize}
\end{enumerate}

\vspace{1em}

\section{Investment Prioritization Framework}

\subsection{Recommended Investment Tiers}

\textbf{Tier 1: Transit Connectivity Priority (502 roads)}
\begin{itemize}[leftmargin=*]
    \item \textbf{Target:} All TOD roads with no sidewalk coverage
    \item \textbf{Impact:} Maximize pedestrian access to Metro-North stations
    \item \textbf{Rationale:} TOD investment requires pedestrian infrastructure
    \item \textbf{Timeline:} 5-7 year capital improvement plan
\end{itemize}

\textbf{Tier 2: Completion Opportunities (352 roads)}
\begin{itemize}[leftmargin=*]
    \item \textbf{Target:} TOD roads with one-side coverage
    \item \textbf{Impact:} Complete pedestrian networks, improve safety
    \item \textbf{Rationale:} Lower cost per mile (infrastructure already half-built)
    \item \textbf{Timeline:} 3-5 year implementation (quick wins)
\end{itemize}

\textbf{Tier 3: Equity \& Non-TOD Expansion}
\begin{itemize}[leftmargin=*]
    \item \textbf{Target:} Non-TOD areas with high pedestrian demand
    \item \textbf{Impact:} Address 81.8\% no-coverage rate in non-TOD zones
    \item \textbf{Rationale:} Environmental justice, equitable access
    \item \textbf{Timeline:} 10-15 year comprehensive plan
\end{itemize}

\vspace{1em}

\section{Key Planning Questions to Explore}

Use the Excel files and PDF reports to investigate:

\begin{enumerate}[leftmargin=*]
    \item \textbf{Spatial Prioritization:} Which TOD roads should be addressed first?
    \begin{itemize}
        \item Filter Excel by coverage status + road type
        \item Overlay with traffic volume data (if available)
        \item Identify clusters of no-coverage areas near high-ridership stations
    \end{itemize}

    \item \textbf{Cost-Benefit Analysis:} What's the ROI of Tier 2 completions vs. Tier 1 new construction?
    \begin{itemize}
        \item Estimate cost per linear foot by road type (Excel file 3)
        \item Calculate cost to complete 352 partial-coverage roads vs. 502 no-coverage roads
        \item Consider pedestrian volume, transit ridership gains
    \end{itemize}

    \item \textbf{Timeline \& Phasing:} What's realistic given budget constraints?
    \begin{itemize}
        \item Scenario A: Focus 100\% on Tier 1 → 5-7 years
        \item Scenario B: Focus 100\% on Tier 2 → 3-5 years (quick wins)
        \item Scenario C: Balanced approach → 10 years comprehensive
    \end{itemize}

    \item \textbf{Grant Eligibility:} Which projects qualify for federal/state funding?
    \begin{itemize}
        \item TOD + transit connectivity grants
        \item ADA compliance improvements
        \item Environmental justice investments (non-TOD equity)
    \end{itemize}

    \item \textbf{Benchmarking:} How do we compare to peer counties?
    \begin{itemize}
        \item See Technical Analysis PDF for comparative data
        \item Industry best practice: 75-85\% coverage in TOD areas
        \item Current Westchester: 54.9\% (gap of 20-30 percentage points)
    \end{itemize}
\end{enumerate}

\vspace{1em}

\section{Methodology Summary}

\textbf{Analysis Method:} DVRPC (Delaware Valley Regional Planning Commission) sidewalk-to-road ratio

\textbf{Coverage Classification:}
\begin{itemize}[leftmargin=*]
    \item \textbf{No Coverage:} Sidewalk-to-road ratio $<$ 0.1 on both sides
    \item \textbf{One-Side:} Ratio 0.4--0.8 on one side, $<$ 0.1 on other
    \item \textbf{Both-Sides:} Ratio $>$ 1.2 on both sides
\end{itemize}

\textbf{TOD Standard:} 0.5 miles (2,640 feet) from Metro-North stations - industry standard for transit-oriented development analysis

\textbf{Data Quality:} 100\% of county road polygons analyzed with geometric validation (EPSG:2260 NY State Plane)

\textbf{Coordinate System:} EPSG:2260 (NY State Plane Long Island, feet) for analysis

\vspace{1em}

\section{Comparative Benchmarks}

\begin{table}[h]
\centering
\begin{tabular}{@{}lrr@{}}
\toprule
\textbf{Area Type} & \textbf{Coverage \%} & \textbf{Gap to Best Practice} \\
\midrule
Westchester TOD Areas & 54.9\% & -20 to -30 percentage points \\
Industry Best Practice (TOD) & 75-85\% & Target range \\
Westchester Non-TOD & 18.2\% & Significant equity gap \\
County-Wide Average & 26.0\% & Below national average \\
\bottomrule
\end{tabular}
\end{table}

\textbf{Key Takeaway:} While Westchester TOD areas perform better than the county average, there's a 20-30 percentage point gap to industry best practices for transit-oriented pedestrian infrastructure.

\vspace{1em}

\section{Next Steps for City Planning}

\begin{enumerate}[leftmargin=*]
    \item \textbf{Review Excel Files:} Explore data, filter by priority criteria
    \item \textbf{Read PDF Reports:} Understand methodology, review recommendations
    \item \textbf{Scenario Planning:} Develop 3-5 investment scenarios with costs
    \item \textbf{Stakeholder Engagement:} Present findings to City Council, community boards
    \item \textbf{Grant Applications:} Identify federal/state funding opportunities
    \item \textbf{Capital Planning:} Integrate into 5-10 year CIP based on priorities
\end{enumerate}

\vspace{2em}

\begin{center}
\rule{0.8\textwidth}{0.4pt}\\
\textbf{Analysis Status: COMPLETE}\\
All 4,386 county roads analyzed with 100\% data coverage\\
October 2025 - Arcanum Research Initiative
\end{center}

\end{document}
